% Part: lambda-calculus
% Chapter: introduction
% Section: term-revisited

\documentclass[../../../include/open-logic-section]{subfiles}

\begin{document}

\olfileid[tr]{lam}{syn}{tr}
\olsection{$\alpha$-Denklik Sınıfları Olarak Terimler}

Bundan sonra terimleri $\alpha$-denkliğine göre ele alacağız. Yani bir terim
yazdığımızda, içinde bulunduğu $\alpha$-denklik sınıfını kastedeceğiz.
Örneğin $\lambd[a][\lambd[b][a c]]$ ile
$\lambd[a][\lambd[b][a c]]$, $\lambd[b][\lambd[a][b c]]$ ve benzeri bütün
$\alpha$-denk terimlerin kümesini göstereceğiz.

Önceki bölümlerde $N,Q$ gibi harfler bir terimi gösterirken bundan sonra bir
sınıfı gösterecek; ileride inceleyeceğimiz nesneler terimler yerine bu
sınıflar olacak. $x,y$ gibi harfler ise değişkenleri göstermeyi sürdürecek.

Ayrıca $\rep{M}$ ile $M$ sınıfının gelişigüzel bir !!{element}{}ını;
birden çoğuna gerek duyarsak $\rep{M}[0]$, $\rep{M}[1]$ ve benzerlerini
göstereceğiz.

Anlatımı sadeleştirmek için terimlere ait gösterimleri yeniden kullanırız.
Sınıflar üzerinde şu tanımları veririz:
\begin{defn}
  \begin{enumerate}
  \item $\lambd[x][N]$, $\lambd[x][\rep{N}]$ terimini içeren sınıf olarak
    tanımlanır.
  \item $PQ$, $\rep{P}\rep{Q}$ terimini içeren sınıf olarak tanımlanır.
  \end{enumerate}
\end{defn}

$\alpha$-dönüşümü bağdaşır olduğu için bu tanımların iyi tanımlı olduğunu
görmek zor değildir.

\begin{defn} \ollabel{def:fv}
  Bir $\alpha$-denklik sınıfı $M$'nin \emph{serbest değişkenleri}, yani
  $\FV{M}$, $\FV{\rep{M}}$ olarak tanımlanır.
\end{defn}

\olref[alp]{thm:fv} içinde gösterildiği gibi
$\FV{\rep{M}[0]}=\FV{\rep{M}[1]}$ olduğundan bu tanım iyi tanımlıdır.

Sınıflar içine yerine koyma gösterimini de yeniden kullanırız:
\begin{defn} \ollabel{def:sub}
  $M$ içindeki $y$ yerine $R$ koyma, yani $\Subst{M}{R}{y}$, yerine koymayı
  tanımlı kılan herhangi bir $\rep{M}$ ve $\rep{R}$ için
  $\Subst{\rep{M}}{\rep{R}}{y}$ olarak tanımlanır.
\end{defn}

\olref[alp]{cor:sub} içinde gösterildiği gibi bu tanım da iyi tanımlıdır.

Bu tanımın akıl yürütmemizi ne kadar sadeleştirdiğine dikkat edin. Örneğin:
\begin{align}
  \Subst{\lambd[x][x]}{y}{x} & =\ollabel{eq:1}\\
  &= \Subst{\lambd[z][z]}{y}{x} \ollabel{eq:2}\\
  &= \lambd[z][\Subst{z}{y}{x}] \\
  &= \lambd[z][z].
\end{align}

\olref{eq:1}, hâlâ terimler üzerinde yerine koyma olarak ele alınırsa
tanımsızdır. Ancak artık sınıflar üzerinde yerine koyma olarak ele alıyoruz;
\olref{eq:2} adımı bu nedenle mümkündür: $\lambd[x][x]$ ile
$\lambd[z][z]$ aynı sınıfa ait olduğundan birini diğeriyle değiştirebiliriz.

Aynı nedenle bundan sonra seçtiğimiz temsilcilerin yerine koyma için gereken
koşulları her zaman sağladığını varsayacağız. Örneğin
$\Subst{\lambd[x][N]}{R}{y}$ gördüğümüzde $\lambd[x][N]$ temsilcisinin
$x\neq y$ ve $x\notin\FV{R}$ olacak biçimde seçildiğini varsayacağız.

$\lambd[x][x]$ nesnesine ``sınıf'' demek biraz tuhaf olduğundan, alışık
olduğumuz $\lambda$-terimlerinden ayırmak üzere bundan sonra bu nesnelere
$\Lambda$-terimleri, kısmın geri kalanında ise yalnızca ``terimler'' diyelim.

\begin{editorial}
  Terimlerle henüz bütünüyle vedalaşamayız: $\Lambda$-terimlerinin bütün
  tanımı $\lambda$-terimlerine dayanır. Ayrıca $\Lambda$-terimleri üzerinde
  fonksiyon tanımlama yöntemi vermedik. Bu nedenle böyle fonksiyonlar önce
  $\lambda$-terimleri üzerinde tanımlanmalı, sonra yerine koymada yaptığımız
  gibi $\Lambda$-terimlerine ``izdüşürülmelidir.'' Bununla birlikte okurun
  $\Lambda$-terimleri üzerinde fonksiyonları nasıl tanımlayabileceğimizi
  sezgisel olarak anlayabileceğini varsayıyoruz.
\end{editorial}
\end{document}
